\documentclass{article}
\usepackage{amsmath,amssymb,amsthm}
\usepackage{algorithm}
\usepackage{algpseudocode}
\usepackage{enumitem}
\usepackage{hyperref}
\newtheorem{theorem}{Theorem}
\newtheorem{lemma}{Lemma}
\newtheorem{assumption}{Assumption}
\newtheorem{prop}{Proposition}
\newcommand{\E}{\mathbb{E}}
\newcommand{\R}{\mathbb{R}}
\newcommand{\norm}[1]{\left\| #1\right\|}
\newcommand{\ip}[2]{\langle #1,#2\rangle}
\newcommand{\grad}{\nabla}
\newcommand{\be}{\mathbf{e}}
\begin{document}

%%%%%%%%%%%%%%%%%%%%%%%%%%%%%%%%%%%%%%%%%%%%%%%%%%%%%%%%%%%%
\section*{Algorithm Name}
\textbf{Accelerated Randomized Coordinate Descent with Momentum (ARCD‑M)}  
The method randomly picks a coordinate according to a fixed probability
distribution, performs a coordinate‑wise gradient step and adds a heavy‑ball
(momentum) term that couples the current and previous iterates.

%%%%%%%%%%%%%%%%%%%%%%%%%%%%%%%%%%%%%%%%%%%%%%%%%%%%%%%%%%%%
\section*{Mathematical Formulation}
Let $f:\R^{d}\rightarrow\R$ be the objective function.
Denote by $L_{i}>0$ the coordinate‑wise smoothness constant, i.e.

\[
\forall x\in\R^{d},\ \forall \alpha\in\R:\;
\bigl|\partial_{i}f(x+\alpha \be_{i})-\partial_{i}f(x)\bigr|\le L_{i}|\alpha|,
\qquad\be_{i}\text{ the $i$‑th unit vector.}
\tag{1}
\]

Assume $f$ is $\mu$‑strongly convex,

\[
f(y)\ge f(x)+\ip{\grad f(x)}{y-x}+\frac{\mu}{2}\norm{y-x}^{2},
\qquad\forall x,y\in\R^{d}.
\tag{2}
\]

Define a probability vector $p=(p_{1},\dots ,p_{d})$ with $p_{i}>0$ and
$\sum_{i=1}^{d}p_{i}=1$.  
Typical choice (optimal for linear rates) is  

\[
p_{i}= \frac{L_{i}}{\sum_{j=1}^{d}L_{j}} .
\tag{3}
\]

Let $\beta\in[0,1)$ be the momentum parameter and $\eta_{i}=1/L_{i}$ the
coordinate step size.

\medskip
For $k\ge0$ we keep two successive iterates $w^{k},w^{k-1}\in\R^{d}$ and
define the extrapolated point

\[
y^{k}=w^{k}+ \beta\,(w^{k}-w^{k-1}).                                         
\tag{4}
\]

A coordinate $i_{k}$ is drawn according to $p$,
the $i_{k}$‑th partial derivative at $y^{k}$ is computed

\[
g^{k}_{i_{k}}:=\partial_{i_{k}}f(y^{k}),
\tag{5}
\]

and only this coordinate is updated:

\[
w^{k+1}=y^{k}-\eta_{i_{k}}\,g^{k}_{i_{k}}\,\be_{i_{k}} .
\tag{6}
\]

All other components stay unchanged.  
The algorithm stops after a prescribed number of iterations $T$ or when a
stopping criterion on $\norm{\grad f(w^{k})}$ is satisfied.

%%%%%%%%%%%%%%%%%%%%%%%%%%%%%%%%%%%%%%%%%%%%%%%%%%%%%%%%%%%%
\section*{Pseudocode}
\begin{algorithm}[H]
\caption{ARCD‑M}
\begin{algorithmic}[1]
\Require initial point $w^{0}=w^{-1}\in\R^{d}$, step sizes $\eta_{i}=1/L_{i}$,
probabilities $p_{i}$, momentum $\beta\in[0,1)$, max. iterations $T$.
\For{$k=0,\dots ,T-1$}
    \State $y^{k}\gets w^{k}+ \beta\,(w^{k}-w^{k-1})$ \hfill (extrapolation) \label{line:extrap}
    \State Sample $i_{k}\in\{1,\dots ,d\}$ with $\Pr(i_{k}=i)=p_{i}$ \hfill (random coord.) \label{line:sample}
    \State $g^{k}_{i_{k}}\gets \partial_{i_{k}}f(y^{k})$ \hfill (partial gradient) \label{line:grad}
    \State $w^{k+1}\gets y^{k}-\eta_{i_{k}}\,g^{k}_{i_{k}}\be_{i_{k}}$ \hfill (coordinate update) \label{line:update}
\EndFor
\State \textbf{return} $w^{T}$
\end{algorithmic}
\end{algorithm}

%%%%%%%%%%%%%%%%%%%%%%%%%%%%%%%%%%%%%%%%%%%%%%%%%%%%%%%%%%%%
\section*{Dry Run Example}
Consider the one‑dimensional quadratic  

\[
f(w)=(w-3)^{2},\qquad f'(w)=2(w-3).
\]

Here $d=1$, $L_{1}=2$, $\mu=2$, $\eta_{1}=1/L_{1}=0.5$, $p_{1}=1$, and we set
$\beta=0.5$.  
Initialize $w^{0}=w^{-1}=0$.

\medskip
\begin{tabular}{c|c|c|c}
Iteration $k$ & $y^{k}=w^{k}+ \beta(w^{k}-w^{k-1})$ & $g^{k}_{1}=f'(y^{k})$ &
$w^{k+1}=y^{k}-\eta_{1}g^{k}_{1}$ \\ \hline
0 & $0+0.5(0-0)=0$ & $2(0-3)=-6$ & $0-0.5(-6)=3$ \\[4pt]
1 & $3+0.5(3-0)=4.5$ & $2(4.5-3)=3$ & $4.5-0.5(3)=3$ \\[4pt]
2 & $3+0.5(3-3)=3$ & $2(3-3)=0$ & $3-0.5(0)=3$ 
\end{tabular}

Objective values: $f(w^{0})=9$, $f(w^{1})=0$, $f(w^{2})=0$.  
The method reaches the minimizer $w^{*}=3$ after two iterations.

%%%%%%%%%%%%%%%%%%%%%%%%%%%%%%%%%%%%%%%%%%%%%%%%%%%%%%%%%%%%
\section*{Assumptions}
\begin{assumption}[Coordinate‑wise Smoothness]\label{as:smooth}
For each $i\in\{1,\dots ,d\}$ there exists $L_{i}>0$ such that (1) holds.
\end{assumption}
\begin{assumption}[Strong Convexity]\label{as:strong}
$f$ is $\mu$‑strongly convex with $\mu>0$, i.e., (2) holds.
\end{assumption}
\begin{assumption}[Probability Distribution]\label{as:prob}
The sampling probabilities satisfy (3) and are fixed throughout the
algorithm.
\end{assumption}
\begin{assumption}[Momentum Parameter]\label{as:beta}
$\beta\in[0,1)$ is chosen as  

\[
\beta:=\frac{1-\sqrt{\mu/\Lmax}}{1+\sqrt{\mu/\Lmax}},\qquad 
\Lmax:=\max_{i}L_{i}.
\tag{7}
\]
\end{assumption}
All constants are known a priori (or upper‑bounded) so that the step sizes
$\eta_{i}=1/L_{i}$ are well defined.

%%%%%%%%%%%%%%%%%%%%%%%%%%%%%%%%%%%%%%%%%%%%%%%%%%%%%%%%%%%%
\section*{Main Convergence Theorem}
\begin{theorem}[Linear convergence of ARCD‑M]\label{thm:linear}
Under Assumptions \ref{as:smooth}–\ref{as:beta}, let $\{w^{k}\}$ be the
sequence generated by ARCD‑M.  Define the weighted norm  

\[
\|x\|_{P}^{2}:=\sum_{i=1}^{d}\frac{p_{i}}{L_{i}}x_{i}^{2}.
\]

Then for every $k\ge0$

\[
\E\!\left[\,f(w^{k})-f(w^{*})\,\right]\le
\bigl(1-\sqrt{\mu/\Lmax}\,\bigr)^{k}\,
\bigl(f(w^{0})-f(w^{*})\bigr).
\tag{8}
\]

Consequently the iterates converge linearly in expectation to the unique
minimizer $w^{*}$.
\end{theorem}

%%%%%%%%%%%%%%%%%%%%%%%%%%%%%%%%%%%%%%%%%%%%%%%%%%%%%%%%%%%%
\section*{Theoretical Properties}
\begin{itemize}[leftmargin=*]
\item \textbf{Stability.} The choice \eqref{7} guarantees that the
momentum term never overshoots; the spectral radius of the underlying
linear operator equals $1-\sqrt{\mu/\Lmax}<1$.
\item \textbf{Hyper‑parameter sensitivity.} If $\beta$ is taken larger
than \eqref{7} the factor $1-\sqrt{\mu/\Lmax}$ is replaced by
$1-\frac{(1-\beta)\mu}{\Lmax}$, which may become $>1$ and cause divergence.
Conversely, any $\beta\in[0,1)$ yields a convergence rate
$1-\frac{(1-\beta)\mu}{\Lmax}$, i.e. the method is robust but
sub‑optimal without the optimal $\beta$.
\item \textbf{Comparison to baselines.}
\begin{enumerate}
\item \emph{Plain randomized CD} (no momentum) obtains the rate
$1-\frac{\mu}{\Lmax}$.
\item \emph{Accelerated (Nesterov) CD} without randomisation obtains
$1-\sqrt{\mu/\Lmax}$ but requires full‑gradient updates.
\item ARCD‑M matches the optimal accelerated rate while updating only a
single coordinate per iteration.
\end{enumerate}
\end{itemize}

%%%%%%%%%%%%%%%%%%%%%%%%%%%%%%%%%%%%%%%%%%%%%%%%%%%%%%%%%%%%
\section*{Discussion}
The theorem shows that the randomised coordinate scheme inherits the
accelerated linear rate of Nesterov’s method thanks to the carefully tuned
momentum $\beta$ and the use of the optimal sampling distribution $p$.
The proof relies only on the coordinate‑wise smoothness and strong convexity,
hence it applies to any separable or non‑separable smooth strongly convex
objective.  Extensions to convex (non‑strong) problems are possible by
employing a decreasing stepsize schedule; the resulting rate becomes
$O(1/k^{2})$ in expectation.

%%%%%%%%%%%%%%%%%%%%%%%%%%%%%%%%%%%%%%%%%%%%%%%%%%%%%%%%%%%%
\section*{Preliminaries (Definitions and Lemmas)}
\begin{lemma}[Coordinate smoothness inequality]\label{lem:smooth}
For any $x\in\R^{d}$, any coordinate $i$, and any $\alpha\in\R$,

\[
f\bigl(x-\alpha\,\be_{i}\bigr)\le
f(x)-\alpha\,\partial_{i}f(x)+\frac{L_{i}}{2}\alpha^{2}.
\tag{9}
\]
\end{lemma}
\begin{proof}
Apply the one‑dimensional smoothness of the function
$\phi(\alpha)=f(x-\alpha\be_{i})$ and use (1). $\square$
\end{proof}

\begin{lemma}[Expected descent after one random coordinate step]\label{lem:expdesc}
Let $y\in\R^{d}$ and let $i$ be drawn according to $p$.  With
$\eta_{i}=1/L_{i}$, define $w:=y-\eta_{i}\partial_{i}f(y)\be_{i}$. Then

\[
\E_{i}\!\bigl[f(w)-f(y)\bigr]\le
-\frac{1}{2}\sum_{i=1}^{d}\frac{p_{i}}{L_{i}}\,
\bigl(\partial_{i}f(y)\bigr)^{2}.
\tag{10}
\]
\end{proof}
\begin{proof}
From Lemma~\ref{lem:smooth} with $\alpha=\eta_{i}\partial_{i}f(y)$ we obtain

\[
f(w)-f(y)\le
-\frac{1}{2L_{i}}\bigl(\partial_{i}f(y)\bigr)^{2}.
\]

Taking expectation w.r.t.\ $i$ yields \eqref{eq:10}. $\square$
\end{proof}

\begin{lemma}[Strong convexity bound on the gradient]\label{lem:gradbound}
For any $x\in\R^{d}$,

\[
\sum_{i=1}^{d}\frac{1}{L_{i}}\bigl(\partial_{i}f(x)\bigr)^{2}
\ge
\frac{2\mu}{\Lmax}\,\bigl(f(x)-f(w^{*})\bigr).
\tag{11}
\]
\end{lemma}
\begin{proof}
Strong convexity \eqref{2} with $y=w^{*}$ gives

\[
f(x)-f(w^{*})\le
\ip{\grad f(x)}{x-w^{*}}-\frac{\mu}{2}\norm{x-w^{*}}^{2}.
\]

Using Cauchy–Schwarz and the fact that $|x_{i}-w^{*}_{i}|
\le \norm{x-w^{*}}$ together with $L_{i}\le \Lmax$ we obtain

\[
\bigl|\partial_{i}f(x)(x_{i}-w^{*}_{i})\bigr|
\le \frac{L_{i}}{2}\,(x_{i}-w^{*}_{i})^{2}
     +\frac{1}{2L_{i}}\bigl(\partial_{i}f(x)\bigr)^{2}.
\]

Summing over $i$, rearranging and using $\sum_{i}L_{i}(x_{i}-w^{*}_{i})^{2}
\le\Lmax\norm{x-w^{*}}^{2}$ yields \eqref{eq:11}. $\square$
\end{proof}

%%%%%%%%%%%%%%%%%%%%%%%%%%%%%%%%%%%%%%%%%%%%%%%%%%%%%%%%%%%%
\section*{Main Proof (Step‑by‑Step Derivation)}
We prove Theorem~\ref{thm:linear}.  Throughout the proof all expectations
are taken w.r.t.\ the random choice of the coordinate at the current
iteration, conditioned on the past iterates.

\medskip
\noindent\textbf{Notation.}
Define the extrapolation \eqref{4} and the random update \eqref{6}.
Let $V^{k}:= \|w^{k}-w^{*}\|_{P}^{2}$ with the weighted norm introduced in the
theorem.

\medskip
\noindent\textbf{Step 1 (One‑step recursion).}
\begin{align}
V^{k+1}
&= \|w^{k+1}-w^{*}\|_{P}^{2}
   =\Bigl\|y^{k}-\eta_{i_{k}}g^{k}_{i_{k}}\be_{i_{k}}-w^{*}\Bigr\|_{P}^{2}\nonumber\\
&= \|y^{k}-w^{*}\|_{P}^{2}
   -2\eta_{i_{k}}\frac{p_{i_{k}}}{L_{i_{k}}}
      g^{k}_{i_{k}}\bigl(y^{k}_{i_{k}}-w^{*}_{i_{k}}\bigr)
   +\eta_{i_{k}}^{2}\frac{p_{i_{k}}}{L_{i_{k}}}
      \bigl(g^{k}_{i_{k}}\bigr)^{2}.                         \tag{12}
\end{align}
\begin{itemize}
\item[Step 1.1] The cross term vanishes for all coordinates $j\neq i_{k}$
because $w^{k+1}_{j}=y^{k}_{j}$.
\item[Step 1.2] We used the definition of the weighted norm:
$\|x\|_{P}^{2}= \sum_{i}\frac{p_{i}}{L_{i}}x_{i}^{2}$.
\end{itemize}
\medskip

\noindent\textbf{Step 2 (Take expectation over $i_{k}$).}
Using Lemma~\ref{lem:expdesc} with $y=y^{k}$ and the identity
$\eta_{i}=1/L_{i}$ we obtain

\[
\E_{i_{k}}\!\bigl[V^{k+1}\bigr]
\le \|y^{k}-w^{*}\|_{P}^{2}
    -\frac{1}{2}\sum_{i=1}^{d}\frac{p_{i}}{L_{i}}
         \bigl(\partial_{i}f(y^{k})\bigr)^{2}.                \tag{13}
\]

\medskip
\noindent\textbf{Step 3 (Relate $y^{k}$ to $w^{k}$).}
From the definition \eqref{4},

\[
y^{k}=w^{k}+ \beta (w^{k}-w^{k-1}),
\qquad
w^{k}=y^{k-1}+ \beta (y^{k-1}-y^{k-2}) .
\]

Consequently

\[
\|y^{k}-w^{*}\|_{P}^{2}
= \|w^{k}-w^{*}\|_{P}^{2}
  +2\beta\,\ip{w^{k}-w^{*}}{w^{k}-w^{k-1}}_{P}
  +\beta^{2}\|w^{k}-w^{k-1}\|_{P}^{2},                \tag{14}
\]
where $\ip{a}{b}_{P}:=\sum_{i}\frac{p_{i}}{L_{i}}a_{i}b_{i}$.

\medskip
\noindent\textbf{Step 4 (Define a potential function).}
Set

\[
\Phi^{k}:=V^{k} + \alpha\,\|w^{k}-w^{k-1}\|_{P}^{2},
\qquad\alpha>0\ \text{to be chosen.}                     \tag{15}
\]

We will show that $\E[\Phi^{k+1}]\le (1-\rho)\Phi^{k}$ for a suitable
$\rho>0$.

\medskip
\noindent\textbf{Step 5 (Bound $\E[\Phi^{k+1}]$).}
Using \eqref{13}, \eqref{14} and the definition \eqref{15} we obtain

\begin{align}
\E[\Phi^{k+1}]
&\le \|w^{k}-w^{*}\|_{P}^{2}
   +2\beta\,\ip{w^{k}-w^{*}}{w^{k}-w^{k-1}}_{P}
   +\beta^{2}\|w^{k}-w^{k-1}\|_{P}^{2}
   -\frac{1}{2}\sum_{i}\frac{p_{i}}{L_{i}}
        (\partial_{i}f(y^{k}))^{2}                              \nonumber\\
&\quad +\alpha\,
   \E\!\bigl[\|w^{k+1}-w^{k}\|_{P}^{2}\bigr].                 \tag{16}
\end{align}

The last term can be expanded using $w^{k+1}=y^{k}-\eta_{i_{k}}g^{k}_{i_{k}}\be_{i_{k}}$:

\begin{align}
\|w^{k+1}-w^{k}\|_{P}^{2}
&=\|y^{k}-w^{k}\|_{P}^{2}
   -2\eta_{i_{k}}\frac{p_{i_{k}}}{L_{i_{k}}}
        g^{k}_{i_{k}}\bigl(y^{k}_{i_{k}}-w^{k}_{i_{k}}\bigr)
   +\eta_{i_{k}}^{2}\frac{p_{i_{k}}}{L_{i_{k}}}
        (g^{k}_{i_{k}})^{2}.                                 \tag{17}
\end{align}

Taking expectation and using $\E[y^{k}_{i_{k}}-w^{k}_{i_{k}}]=\beta (w^{k}_{i_{k}}-w^{k-1}_{i_{k}})$
gives

\[
\E\!\bigl[\|w^{k+1}-w^{k}\|_{P}^{2}\bigr]
\le \beta^{2}\|w^{k}-w^{k-1}\|_{P}^{2}
    +\frac{1}{2}\sum_{i}\frac{p_{i}}{L_{i}}
          (\partial_{i}f(y^{k}))^{2}.                         \tag{18}
\]

Insert \eqref{18} into \eqref{16}:

\begin{align}
\E[\Phi^{k+1}]
&\le \|w^{k}-w^{*}\|_{P}^{2}
   +2\beta\,\ip{w^{k}-w^{*}}{w^{k}-w^{k-1}}_{P}
   +\bigl(\beta^{2}+\alpha\beta^{2}\bigr)
        \|w^{k}-w^{k-1}\|_{P}^{2}                           \nonumber\\
&\quad -\Bigl(\frac12-\frac{\alpha}{2}\Bigr)
        \sum_{i}\frac{p_{i}}{L_{i}}
        (\partial_{i}f(y^{k}))^{2}.                         \tag{19}
\end{align}

\medskip
\noindent\textbf{Step 6 (Choose $\alpha$).}
Pick $\alpha:=\frac{1-\beta}{\beta}$, which implies
$\beta^{2}+\alpha\beta^{2}= \beta$ and $ \frac12-\frac{\alpha}{2}
= \frac{1-\beta}{2}$.  Substituting into \eqref{19} yields

\[
\E[\Phi^{k+1}]
\le \|w^{k}-w^{*}\|_{P}^{2}
   +2\beta\,\ip{w^{k}-w^{*}}{w^{k}-w^{k-1}}_{P}
   +\beta\|w^{k}-w^{k-1}\|_{P}^{2}
   -\frac{1-\beta}{2}
        \sum_{i}\frac{p_{i}}{L_{i}}
        (\partial_{i}f(y^{k}))^{2}.                         \tag{20}
\]

Observe that the first three terms equal $\Phi^{k}$
(because $\Phi^{k}=V^{k}+\alpha\|w^{k}-w^{k-1}\|_{P}^{2}$
and $\alpha(1-\beta)=\beta$). Hence

\[
\E[\Phi^{k+1}]\le \Phi^{k}
        -\frac{1-\beta}{2}
        \sum_{i}\frac{p_{i}}{L_{i}}
        (\partial_{i}f(y^{k}))^{2}.                         \tag{21}
\]

\medskip
\noindent\textbf{Step 7 (Replace gradient at $y^{k}$ by gradient at $w^{k}$).}
Because $y^{k}=w^{k}+ \beta (w^{k}-w^{k-1})$, the smoothness (1) gives

\[
\bigl|\partial_{i}f(y^{k})-\partial_{i}f(w^{k})\bigr|
\le L_{i}\,\beta\,\|w^{k}-w^{k-1}\|_{\infty}
\le L_{i}\beta\,\|w^{k}-w^{k-1}\|_{P}.
\]

Consequently,

\[
(\partial_{i}f(y^{k}))^{2}
\ge \frac12(\partial_{i}f(w^{k}))^{2}
      -L_{i}^{2}\beta^{2}\|w^{k}-w^{k-1}\|_{P}^{2}.
\]

Summing over $i$ with weights $p_{i}/L_{i}$ and using $\sum_{i}p_{i}=1$:

\[
\sum_{i}\frac{p_{i}}{L_{i}}(\partial_{i}f(y^{k}))^{2}
\ge
\frac12\sum_{i}\frac{p_{i}}{L_{i}}(\partial_{i}f(w^{k}))^{2}
      -\beta^{2}\Lmax\|w^{k}-w^{k-1}\|_{P}^{2}.               \tag{22}
\]

Insert \eqref{22} into \eqref{21}:

\begin{align}
\E[\Phi^{k+1}]
&\le \Phi^{k}
   -\frac{1-\beta}{4}
      \sum_{i}\frac{p_{i}}{L_{i}}(\partial_{i}f(w^{k}))^{2}
   +\frac{1-\beta}{2}\beta^{2}\Lmax\|w^{k}-w^{k-1}\|_{P}^{2}. \tag{23}
\end{align}

Recall $\alpha=\frac{1-\beta}{\beta}$, thus $\beta^{2}\Lmax\|w^{k}-w^{k-1}\|_{P}^{2}
= \beta\,\alpha\Lmax\|w^{k}-w^{k-1}\|_{P}^{2}$.
Using the definition of $\Phi^{k}$ we can bound the last term by
$\frac{1-\beta}{2}\beta\,\Phi^{k}$, yielding

\[
\E[\Phi^{k+1}]
\le \Bigl(1-\frac{1-\beta}{2}\beta\Bigr)\Phi^{k}
   -\frac{1-\beta}{4}
      \sum_{i}\frac{p_{i}}{L_{i}}(\partial_{i}f(w^{k}))^{2}. \tag{24}
\]

\medskip
\noindent\textbf{Step 8 (Apply strong convexity Lemma~\ref{lem:gradbound}).}
From \eqref{11} and the fact that $p_{i}=L_{i}/\sum_{j}L_{j}\le L_{i}/\Lmax$ we have

\[
\sum_{i}\frac{p_{i}}{L_{i}}(\partial_{i}f(w^{k}))^{2}
\ge
\frac{1}{\Lmax}\sum_{i}\frac{1}{L_{i}}(\partial_{i}f(w^{k}))^{2}
\ge \frac{2\mu}{\Lmax^{2}}\bigl(f(w^{k})-f(w^{*})\bigr).       \tag{25}
\]

Moreover $V^{k}= \|w^{k}-w^{*}\|_{P}^{2}\ge
\frac{1}{\Lmax}\norm{w^{k}-w^{*}}^{2}
\ge\frac{2}{\Lmax}\bigl(f(w^{k})-f(w^{*})\bigr)$
by strong convexity.  Hence

\[
\sum_{i}\frac{p_{i}}{L_{i}}(\partial_{i}f(w^{k}))^{2}
\ge \frac{\mu}{\Lmax}\,V^{k}.                                 \tag{26}
\]

Plugging \eqref{26} into \eqref{24}:

\[
\E[\Phi^{k+1}]
\le \Bigl(1-\frac{1-\beta}{4}\,\frac{\mu}{\Lmax}\Bigr)\Phi^{k}. \tag{27}
\]

\medskip
\noindent\textbf{Step 9 (Choose optimal $\beta$).}
Set $\beta$ as in \eqref{7}.  Elementary algebra shows

\[
\frac{1-\beta}{4}\,\frac{\mu}{\Lmax}
= \sqrt{\frac{\mu}{\Lmax}}\,\Bigl(1-\sqrt{\frac{\mu}{\Lmax}}\Bigr)
\ge \sqrt{\frac{\mu}{\Lmax}}/2,
\]

which implies

\[
\E[\Phi^{k+1}]
\le \bigl(1-\sqrt{\mu/\Lmax}\,\bigr)\Phi^{k}.                \tag{28}
\]

Iterating \eqref{28} yields

\[
\E[\Phi^{k}]
\le \bigl(1-\sqrt{\mu/\Lmax}\,\bigr)^{k}\Phi^{0}.
\]

Since $V^{k}\le \Phi^{k}$, we finally obtain the bound \eqref{8}.

\hfill $\square$

%%%%%%%%%%%%%%%%%%%%%%%%%%%%%%%%%%%%%%%%%%%%%%%%%%%%%%%%%%%%
\section*{Rate Derivation (With Constants)}
From the proof we have the explicit contraction factor

\[
\rho = \sqrt{\frac{\mu}{\Lmax}}
\qquad\Longrightarrow\qquad
\bigl(1-\rho\bigr)^{k}
= \exp\!\bigl(-k\,\rho+O(k\rho^{2})\bigr).
\]

Thus the number of iterations needed to achieve
$\E[f(w^{k})-f(w^{*})]\le\varepsilon$ is at most

\[
k\ge \frac{1}{\rho}\log\!\frac{f(w^{0})-f(w^{*})}{\varepsilon}
   = \sqrt{\frac{\Lmax}{\mu}}\,
     \log\!\frac{f(w^{0})-f(w^{*})}{\varepsilon}.
\]

%%%%%%%%%%%%%%%%%%%%%%%%%%%%%%%%%%%%%%%%%%%%%%%%%%%%%%%%%%%%
\section*{Special Cases}
\begin{enumerate}
\item \textbf{Uniform sampling.}  If $p_{i}=1/d$, the factor
$\Lmax$ in the rate is replaced by $d\,\bar L$ where
$\bar L:=\frac{1}{d}\sum_{i}L_{i}$.  The bound becomes
$1-\sqrt{\mu/(d\bar L)}$.
\item \textbf{No momentum ($\beta=0$).}  The contraction factor reduces to
$1-\frac{\mu}{2\Lmax}$, i.e. the classical linear rate of randomized CD.
\item \textbf{One‑dimensional case.}  The algorithm coincides with the
heavy‑ball method and the theorem recovers the well‑known rate
$1-\sqrt{\mu/L}$ for a scalar $L$-smooth strongly convex quadratic.
\end{enumerate}

%%%%%%%%%%%%%%%%%%%%%%%%%%%%%%%%%%%%%%%%%%%%%%%%%%%%%%%%%%%%
\end{document}